% SPDX-License-Identifier: Apache-2.0
% covers: CfgFlash, FlashPins, Startup
\datasheet{CfgFlash}{The configuration flash reached from user logic after the bitstream has loaded, through \code{STARTUPE2}.}

\dsfacts{\code{lib/parts/src/cfgflash.rs}}{\code{txhdl\_parts::cfgflash::CfgFlash}}%
{\code{//docs:examples}}

\subsection*{Function}

On a 7-series part the configuration flash's clock is the dedicated
CCLK pin, which user logic reaches only through the \code{STARTUPE2}
primitive's \code{USRCCLKO} (issue 312). The primitive brings two
rules, from UG470: its \code{EOS} says when configuration is over, and
until then the flash is not the user's; and the first three clocks on
\code{USRCCLKO} after \code{EOS} never reach the pin, so a command whose
clocks start there reaches the chip three bits short.

\code{CfgFlash} answers both, in two children joined as a unit of
units. \code{prim} is \code{Startup}, the primitive as the foreign
module \code{startup} in \code{lib/parts/hdl/startup.v}: synthesised,
it instantiates \code{STARTUPE2} and registers \code{EOS} twice;
simulated, in Rust, Verilog and VHDL alike, it is a model that raises
\code{eos} after \code{EOS\_AFTER} cycles, swallows the three clocks, and
brings out the pin as \code{cclk} for a model of the chip. \code{pins}
is \code{FlashPins<GUARD>}, between the chip and the two parts that
reach it, the SPI master \code{Spi} and the window \code{FlashWin}.

\subsection*{Parameters}

\begin{center}\footnotesize
\begin{tabular}{@{}lp{0.7\textwidth}@{}}
\toprule
Parameter & Meaning \\
\midrule
\code{GUARD} & One refuses write enable, \code{0x06}, from the master;
zero passes it. The tables use 1, as the board and the example do. \\
\bottomrule
\end{tabular}
\end{center}

\code{EOS\_AFTER}, eight, and \code{SWALLOWED}, three, are constants of
\code{cfgflash.rs} and belong to the model.

\subsection*{Register map}

None. The master and the window have their own; this part has wires.

\dstables{CfgFlash}

\subsection*{Behaviour}

\begin{itemize}
\item Until \code{eos} the pins hold the chip unselected and the clock
low, and \code{w\_rst} high, so the window takes nothing off its bus.
\item After \code{eos} they turn the clock eight times, four clocks, with
the chip not selected, so the three the part swallows are spent on
nothing. Then \code{w\_rst} falls.
\item From then on the wires go to whichever part selects the chip
first, the window before the master when both ask in one cycle, and
stay with it until it lets go of its select. The other part's select
is not passed on.
\item With \code{GUARD} set, the pins count the master's command a bit
a rise of its clock. On the seventh rise of a command whose seven bits
begin write enable they raise the select, so the chip has a command
shorter than eight bits, which it discards. \code{refused} is high
until the master lets go. Every command that programs or erases a
flash, or writes its registers, is ignored by the chip without write
enable, so that one refusal is the door to all of them.
\item Every output of the pins is a register, and the clock passes the
primitive as well, so the chip sees the master's lines a cycle late
and its clock two. A master must run at a half bit of at least four
cycles, \code{DIV} of three, for the chip's answer to be on \code{miso}
when it takes it. The chip's \code{miso} goes to both parts directly.
\item Only mode 0 is supported: the clock idles low.
\item The unit runs the pins before the primitive, and the primitive
drives \code{eos} a step ahead, since the two drive each other through
wires and a wire promises nothing about which process ran first.
\end{itemize}

\subsection*{Verification}

\begin{itemize}
\item \code{//lib/examples:ex\_cfgflash} puts the master behind its
bridge and the window on its own link, both on one \code{FlashDevice}
through \code{CfgFlash<1>}. A window read asked before \code{eos} waits
and is answered; the master reads the identity \code{20 ba 18} whole,
the first command after the swallowed clocks; its write enable is cut
short and \code{refused} rises once; a second window read is answered.
The chip took the identity and the two reads whole and one command
short.
\item \code{//docs:cfgflash\_sim\_cfgflash\_tb\_test} and
\code{//docs:cfgflash\_vsim\_test} simulate the netlist against that
run under nvc and Verilator, with the startup model in both languages.
\item \code{//cpu/vreteno:board\_test} runs
\code{the\_configuration\_flash\_answers\_on\_the\_board} through
\code{Board}, and \code{//docs:vreteno} reports the same program on the
board: the identity, the latch clear and the bitstream's sync word.
\end{itemize}

\subsection*{Used by}

\code{Board}, as \code{cfgflash}, with \code{spi} and \code{flashwin}
on its other side and the configuration pins on the top.
\code{ex\_cfgflash}.

\subsection*{Limits}

The guard refuses write enable only from the master; the window only
reads. A command that a chip takes without write enable is not
refused. The window and the master share the chip by select, so a
master command while the window holds the chip is not passed on and
reads back what the line holds, and a program that uses the master
must not do so while the window is being read. Synthesised,
\code{cclk} is held low and means nothing; it exists for simulation.
